\documentclass[11pt]{article}
\usepackage[a4paper,margin=24mm]{geometry}
\usepackage{amsmath,amssymb,amsthm,mathtools,longtable,array,booktabs}
\usepackage[hidelinks]{hyperref}
\usepackage[T1]{fontenc}
\newtheorem{theorem}{Theorem}[section]
\newtheorem{proposition}[theorem]{Proposition}
\newtheorem{corollary}[theorem]{Corollary}
\newtheorem{lemma}[theorem]{Lemma}
\theoremstyle{definition}
\newtheorem{example}[theorem]{Counterexample}
\newcommand{\E}{\mathbb E}
\newcommand{\Pp}{\mathbb P}
\newcommand{\1}{\mathbf 1}
\newcommand{\sg}{\operatorname{sgn}_{+}}
\newcommand{\pos}[1]{\left(#1\right)_{+}}
\newcommand{\pr}{\textbf{[PROVED]}}
\newcommand{\ce}{\textbf{[COUNTEREXAMPLE]}}
\newcommand{\nc}{\textbf{[NUMERICALLY CHECKED]}}
\newcommand{\cj}{\textbf{[CONJECTURE]}}
\DeclareMathOperator{\conv}{conv}
\DeclareMathOperator{\Per}{Per}
\title{Marginal information and acquisition\\Rounds 1--2: proof appendix and Week 15 cross-audit}
\author{Research working note; priority not established}
\date{4 October 2026}
\begin{document}
\maketitle
\begin{abstract}
For a finite binary pool, expected Hamming-error reduction is determined by
first moments and uncentered pair moments. This note proves the sharp one-query
margin guarantee, the balanced marginal-only minimax value, an exact formula
for homogeneous rare-class marginals, and misspecification bounds.
The general sharp finite-horizon margin constant remains unresolved.
It audits Week 15 and distinguishes coherent conditioning from Laplace
refitting. This document and all computations are outside the repository;
no empirical acquisition run was repeated.
\end{abstract}
\tableofcontents
\section{Scope, status, and priority audit}
All laws are conditional on the current data unless stated otherwise.
The default model has a fixed pool of $N\ge2$ binary labels, noiseless
revelation of the queried label, equal-weight evaluation on \emph{all} $N$
points including queried points, and ideal conditioning under one joint law.
Each realized world may be deterministic while the unknown label vector is
random epistemically. A spin takes values $\{-1,+1\}$; transform $0/1$ labels
before using spin formulas. Set $\sg(0)=+1$.

Ratios compare expected \emph{risk reductions}, not residual errors.
Zero optimal gain gives a vacuous guarantee. Hamming loss is not NSD, ASSD,
or boundary Dice. Statements about other targets, noise, weights or model
updates explicitly change the default assumptions.
\pr\ means proved here, \ce\ a refuted generalization with a construction,
\nc\ a finite computation, and \cj\ an unresolved claim.
Mathematical correctness does not establish publication priority.

\subsection{Closest results and verdicts}
This is a bounded primary-source audit, not a certificate of absence.
\small
\begin{longtable}{p{.24\textwidth}p{.36\textwidth}p{.31\textwidth}}
\toprule Source & Closest result & Verdict\\\midrule\endhead
Mussmann--Liang, ICML 2018 \cite{ml18icml}&
Theorems 4.1--4.2: asymptotic uncertainty-sampling data efficiency versus
passive logistic regression.&
Different comparator and regime; not a fixed-pool one-step EER ratio.\\
Mussmann--Liang, NeurIPS 2018 \cite{ml18nips}&
Uncertainty sampling as preconditioned SGD on smoothed zero-one loss.&
Optimization/stationarity results, not the $1/(N-1)$ constant.\\
Mussmann et al., 2022 \cite{eer}&
Section 3.2, PDF equation (6), MEZL: expected zero-one loss via pairwise
predictive probabilities.&
Direct priority for pairwise EER. Our binary max identity is a specialization.\\
Liu--Li, 2023 v1 \cite{liu}&
Equivalent losses and excess-risk guarantees for uncertainty weighting.&
Different optimization model; no matching finite-pool minimax theorem located.
Do not attribute later revisions to the original version.\\
Bickford Smith et al., EPIG 2023 \cite{epig}&
Section 4: predictive mutual information targeted to test inputs, and expected
cross-entropy after querying.&
Direct priority for target-aware dependence; not zero-one sharp constants.\\
H\"ubotter et al., 2024 \cite{tal}&
Conference Theorems 3.2--3.3: transductive variance/confidence bounds with
distinct sampling and target sets.&
Structural transfer guarantees for function observations; not binary
Hamming ratios without extra assumptions.\\
Wang--Sun--Grosse, 2021 \cite{wang}&
Predictive-correlation benchmarking through TAL, metacorrelations, and XLL.&
Strong conceptual priority for the failure of marginal quality to imply
acquisition quality; also separates selection from prediction models.\\
Pass--Spektor \cite{pass}&
Extremal moments of pairwise independent Rademacher sums and exchangeable
sign-slice constructions.&
Close extremal method, different objective. Ingredients are not new.\\
Cut-polytope gap inequalities \cite{gap}&
Integer/parity strengthening of positive semidefiniteness.&
The parity step for $\kappa_N$ is a known gap inequality.\\
Bickford Smith--Foster--Rainforth 2024 \cite{update}&
Section 4.2.1: mismatch between assumed Bayesian updates and actual learning.&
Direct conceptual priority for the update-consistency caveat.\\
Hu--Mussmann, arXiv v3, 2026 \cite{batch}&
Myopic Bayesian decision theory and partial batch label sampling for EPIG.&
Further decision-theoretic prior art; not the optimal adaptive Hamming ratio.\\
\bottomrule
\end{longtable}
\normalsize
No equivalent exact $1/(N-1)$ or $\kappa_N$ statement was located in these
sources or targeted searches under EER competitive ratios, cut-matrix norm
ratios, sign-sum extremality, or oblivious/marginal-only selection.
\emph{Priority remains unestablished.} The general claim that marginals miss
acquisition-relevant dependence is already established prior art.
The rare-marginal results below also need a separate priority audit.

\section{Binary identity and sharp one-step margin bound}
\begin{theorem}[Arbitrary binary target and observation]\label{th:binary}
\pr\ Let $T,O\in\{-1,+1\}$ have any joint law, $m=\E T$, $c=\E TO$.
The minimum expected zero-one error after observing $O$ is
\[
 e(T\mid O)=\frac{1-\max\{|m|,|c|\}}2.
\]
For targets $T_i$, weights $w_i\ge0$, $\sum_iw_i=1$, and observations $O_j$,
put $m_i=\E T_i$ and $c_{ij}=\E T_iO_j$. Then
\begin{align}
 R_0&=\tfrac12\sum_iw_i(1-|m_i|),\\
 F(j)&=\tfrac12\sum_iw_i[1-\max\{|m_i|,|c_{ij}|\}],\\
 \Delta(j)&=\tfrac12\sum_iw_i\pos{|c_{ij}|-|m_i|}.\label{eq:gain}
\end{align}
\end{theorem}
\begin{proof}
$\E[T\1_{\{O=s\}}]=(m+sc)/2$. Optimal accuracy minus error therefore equals
$\sum_s|\E[T\1_{\{O=s\}}]|=(|m+c|+|m-c|)/2=\max(|m|,|c|)$.
Zero-probability branches cause no problem. Sum the target risks.
\end{proof}
\begin{theorem}[Sharp margin guarantee]\label{th:margin}
\pr\ In the default model put $u_i=\min(p_i,1-p_i)$ and choose
$a\in\arg\max_i u_i$. Then
$\Delta(a)\ge\max_j\Delta(j)/(N-1)$.
The infimum is sharp even with a unique margin maximizer.
For $N=2$, margin is optimal.
\end{theorem}
\begin{proof}
Let $e_i^{(j)}$ be expected posterior error, $v_{ij}=u_i-e_i^{(j)}$ and
$G_j=\sum_i v_{ij}=N\Delta(j)$.
Then $0\le v_{ij}\le u_i$, $v_{jj}=u_j$, and $u_a\ge u_b$ implies
$e_b^{(a)}\le e_a^{(b)}$ by Theorem \ref{th:binary}.
Thus $u_b+v_{ab}\le u_a+v_{ba}$ and
\[
 G_b\le u_b+v_{ab}+(N-2)u_a
 \le G_a+(N-2)u_a\le(N-1)G_a.
\]
For sharpness take an independent fair decoy $D$ and $N-1$ copies of
$Z\sim\mathrm{Bernoulli}(1/2+\eta)$. Margin uniquely queries $D$.
Gains are $1/(2N)$ and $(N-1)(1/2-\eta)/N$, whose ratio tends to $1/(N-1)$.
\end{proof}
\begin{corollary}[Independence, entropy, perturbations]\label{cor:basic}
\pr\ Pairwise independence makes margin one-step Hamming-optimal.
Noiseless margin maximizes reduction of joint label entropy for every joint
law. If $|\widehat m_i-m_i|\le\epsilon_m$ and
$|\widehat c_{ij}-c_{ij}|\le\epsilon_c$, each score \eqref{eq:gain} has error
at most $(\epsilon_m+\epsilon_c)/2$, and its estimated maximizer has regret
at most $\epsilon_m+\epsilon_c$.
\end{corollary}
\begin{proof}
Under independence $c_{ij}=m_im_j$ for $i\ne j$, so off-diagonal gains vanish.
For entropy, $I(Y;Y_j)=H(Y_j)$, increasing with $u_j$.
For perturbations use the Lipschitz properties of absolute value and positive
part, then compare the estimated and true maximizing scores.
\end{proof}

\section{Balanced minimax and the cut polytope}
\begin{theorem}[Exact balanced marginal-only minimax]\label{th:kappa}
\pr\ Suppose all $p_i=1/2$. A marginal-only selector uses probabilities
$q_j$, without access to the joint law. Put $M=N-1$; let $a$ be the greatest
nonnegative integer at most $\sqrt M$ with the same parity as $M$, and $b=a+2$.
Set
\[
 \tau=(M+ab)/(a+b),\qquad \kappa_N=\frac{N+2\tau}{N(1+\tau)}.
\]
Then uniform selection is minimax and
\[
 \max_{q\in\Delta_N}\inf_P
 \frac{\sum_jq_j\Delta_P(j)}{\max_j\Delta_P(j)}=\kappa_N.
\]
Equivalently, with $\mathcal C_N=\conv\{ss^\top:s\in\{-1,+1\}^N\}$,
\[
 \min_{C\in\mathcal C_N}
 \frac{\sum_{ij}|C_{ij}|}{N\max_k\sum_i|C_{ik}|}=\kappa_N.
\]
\end{theorem}
\begin{proof}
Global sign symmetrization makes means zero without changing second moments,
so $\mathcal C_N$ is exactly the feasible set. Let a maximal column $k$ have
sum $1+t$ and let $A=\sum_{ij}|C_{ij}|$. Diagonal and root entries imply
$A\ge N+2t$. Let $W=\sum_{i\ne k}\sigma_iS_iS_k$ with
$\sigma_iC_{ik}=|C_{ik}|$. Then $\E W=t$ and $W$ lies on the parity-$M$
lattice, so $(W-a)(W-b)\ge0$. Consequently
\[
 \E W^2\ge(a+b)t-ab,\qquad
 A\ge1+2t+\E W^2\ge1+(2+a+b)t-ab.
\]
Divide by $N(1+t)$: the first bound decreases and the second increases.
Their intersection is $t=\tau$, yielding the lower bound.

For attainment take a fair root $R$ and an independent sign vector $V$ of
length $M$. Its sum $W$ takes values $a,b$ with
$\Pp(W=b)=(M-a^2)/(b^2-a^2)$; conditional on $W$, use a uniform sign slice.
Set root spin $R$ and leaf spins $-RV_i$. Then $\E W=\tau$, $\E W^2=M$,
leaf-pair correlations vanish, and root--leaf correlations equal $-\tau/M$.
The root and leaf column sums are $1+\tau$ and $1+\tau/M$.
Uniform selection attains the bound. For any $q$ place the root at an index
with $q_k\le1/N$, making its ratio no larger than the uniform ratio.
Every world in this construction contains both classes.
\end{proof}
\begin{proposition}[PSD relaxation and known gap inequality]\label{prop:psd}
\pr\ Over all PSD unit-diagonal matrices the same minimum is
$(1+\sqrt{N-1})/N$. It equals $\kappa_N$ exactly when $N-1$ is a square.
The parity step above is a cut-polytope gap inequality.
\end{proposition}
\begin{proof}
The Schur complement gives the algebraic analogue $\E W^2\ge t^2$, hence
the ratio is at least
$\max\{(N+2t)/(N(1+t)),(1+t)/N\}$.
These intersect at $t=\sqrt M$. The matrix with leaf block $I_M$ and every
root--leaf entry $1/\sqrt M$ is PSD and attains the bound.
For $N>2$, binary equality requires $W=\sqrt M$ almost surely, possible only
for square $M$; the extremal slice construction then attains it.
The case $N=2$ has value one directly.
For the gap interpretation take integer coefficients $b_k=-(a+1)$,
$b_i=\sigma_i$. The odd integer $b^\top S=S_k(W-a-1)$ satisfies
$(b^\top S)^2\ge1$, precisely $(W-a)(W-a-2)\ge0$.
These inequalities predate this note \cite{gap}.
\end{proof}
\nc\ Round-1 checks of extremal laws for $N=4,\ldots,10$ agree with this proof.
Examples: $\kappa_3=5/6$, $\kappa_4=7/10$, $\kappa_5=3/5$,
$\kappa_6=5/9$, $\kappa_7=1/2$, $\kappa_{10}=2/5$.

\section{General marginals and an exact rare-class formula}
\begin{proposition}[Fully oblivious versus marginal-only selection]
\label{prop:oblivious}
\pr\ A selector fixed before seeing either marginals or dependence has
minimax gain ratio $1/N$ over all binary pool laws. Uniform attains it.
This is a different game from the balanced marginal-only game.
For any gain guarantee $\Delta_a\ge r\Delta^\star$, the corresponding
residual-risk statement is
$R_a\le(1-r)R_0+rR^\star$, not $R_a/R^\star\le1/r$.
\end{proposition}
\begin{proof}
Uniform averages nonnegative gains, hence achieves at least their maximum
divided by $N$. For any fixed selector choose its least-probable coordinate
as an independent fair bit and make every other label deterministic.
Only that query has nonzero value, giving ratio at most $1/N$.
If all marginals must be nondegenerate, approximate the deterministic ones
by independent bits of vanishing uncertainty. The risk conversion follows
by subtracting gains from $R_0$.
\end{proof}
\begin{theorem}[Finite LP characterization]\label{th:LP}
\pr\ Fix $p=(p_1,\ldots,p_N)$ with some nondegenerate coordinate. Let
$\mathcal P(p)$ be all world laws with these marginals. Partition this
probability simplex by $c_{ij}=\pm|m_i|$, and collect the vertices $v$ of
all cells. Set $G_j(v)=N\Delta_v(j)$ and
$r_j(v)=G_j(v)/\max_kG_k(v)$. The minimax value and an optimal selector solve
\[
 \max_{q,z}z,\quad q_j\ge0,\quad\sum_jq_j=1,\quad
 \sum_jq_jr_j(v)\ge z\quad\hbox{for every cell vertex }v.
\]
An optimizer can be constant within each identical-marginal group.
For two levels, selection therefore reduces to one group-mass variable,
although enumerating constraints remains exponential.
\end{theorem}
\begin{proof}
Each $G_j$ is linear on each cell. For fixed $q$, the ratio equals
$\min_{k:G_k>0}(q^\top G)/G_k$. A linear-fractional function with nonnegative
denominator attains its infimum at a positive-denominator vertex: express
any point as a convex combination of vertices and weight their ratios by
their denominators. Zero-denominator vertices add a nonnegative numerator.
At every law some $G_j\ge u_j>0$. Taking all cells proves the LP.
The worst-law value is concave in $q$ and invariant under within-group
permutations; averaging an optimizer proves the symmetry statement.
\end{proof}
\begin{proposition}[Uniform is not always minimax]\label{prop:uniform}
\pr\ If all $p_i=p\in(0,1)$, uniform selection is minimax; the same holds
when all $u_i$ are equal, after known coordinate label flips.
For $N=2$, $u_1\ge u_2>0$, the minimax value is one, attained by margin.
A selector choosing point 1 with probability $q$ has worst-case ratio
$q+(1-q)u_2/u_1$.
\end{proposition}
\begin{proof}
Average over all permutations in the homogeneous case; flips preserve risks.
For $N=2$, margin is optimal. Let $e_i^{(j)}$ be expected posterior errors.
Then $e_1^{(2)}\ge e_2^{(1)}$, and
$G_1-G_2=e_1^{(2)}-e_2^{(1)}\le u_1-u_2$.
The last bound follows directly by subtracting the two max formulas.
Since $G_1\ge u_1$, $G_2/G_1\ge u_2/u_1$.
Independent labels attain equality.
\end{proof}
\ce\ At $(p_1,p_2)=(1/2,1/100)$, uniform has value $0.51$ while margin has
value one. Thus heterogeneous imbalance can make uniform suboptimal.

\begin{theorem}[A nontrivial exact two-level family]\label{th:twolevel}
\pr\ For $N=3$, marginals $(a,a,b)$, $0<2a\le b\le3a$ and $b\le1/3$,
put $r=b/a\in[2,3]$. An optimal marginal-only selector assigns total mass
\[
 \theta=\frac{r(3-r)}{6r-3-r^2}
\]
equally to the first two points and mass $1-\theta$ to the third.
Its exact minimax value is
\[
 V(a,a,b)=\frac{2r}{6r-3-r^2}.
\]
In particular, for $(a,a,b)=(.1,.1,.25)$, the selector is
$(5/46,5/46,18/23)$ and its value is $20/23$.
Pure margin has value $5/6$ in this example; uniform has value $3/5$.
\end{theorem}
\begin{proof}
All probabilities are at most $1/3$, so the influence of query $j$ on
target $i$ is $(2q_{ij}-p_j)_+$, where $q_{ij}=P(Y_i=Y_j=1)$.
Since $b\ge2a$, $G_3=b$ always. Put
$x=(2q_{12}-a)_+$ and $y_i=(2q_{i3}-a)_+$.
Then $G_1=a+x+y_1$, $G_2=a+x+y_2$, with $0\le x,y_i\le a$.
Symmetry permits a selector with masses $(\theta/2,\theta/2,1-\theta)$.
Independent labels give ratio $1-\theta(1-1/r)$.
Taking $Y_1=Y_2$ with their common event a subset of $\{Y_3=1\}$ gives
ratio $\theta+(1-\theta)r/3$.
Balancing these two lines gives the asserted $\theta$ and an upper bound
on every selector's worst-case ratio.

For the matching lower bound, if $\max(G_1,G_2)\le b$, the ratio is at
least the independent-label line since each $G_i\ge a$.
Otherwise suppose $d=G_1>b\ge2a$. Both $x,y_1$ are positive and
$q_{23}\ge q_{12}+q_{13}-a$ implies
$y_2\ge x+y_1-a=d-2a$.
Also $x\ge d-2a$, so $G_2\ge2d-3a$.
The ratio is therefore at least
\[
 \frac{\theta(3d-3a)/2+(1-\theta)b}{d}.
\]
For the specified $\theta$, $0\le\theta\le2/5$ and
$b(1-\theta)-3a\theta/2>0$.
This expression decreases with $d\le3a$, and its value at $3a$ is exactly
the second line. This proves the minimax formula.
For pure margin $b/\max G\ge b/(3a)$, attained by the nested law.
For uniform in the stated example, if the maximum is $b$ the ratio is at
least $(2a+b)/(3b)=3/5$. If it is $d>b$, the preceding bound gives
$1+(b-3a)/(3d)>3/5$. Independence attains $3/5$.
\end{proof}

\begin{theorem}[Exact homogeneous rare-class value]\label{th:rare}
\pr\ Suppose $p_i=p$, $0<p\le1/3$, and denote the minimax value by
$\kappa_N(p)$. For $z\ge0$ define
$\phi(z)=2\lfloor z\rfloor z-\lfloor z\rfloor(\lfloor z\rfloor+1)$.
For $1\le K\le N-1$, $0\le t\le K$, put
\[
 \mu_1=(K+t)/2,\quad
 \mu_0=\frac{p(K-t)}{2(1-p)},\quad
 L_{K,p}(t)=\pos{2p\phi(\mu_1)+2(1-p)\phi(\mu_0)-pK(K-1)}.
\]
Then
\begin{equation}\label{eq:rare}
 \boxed{\kappa_N(p)=
 \min\left\{1,\
 \min_{\substack{1\le K\le N-1,\ 0\le t\le K\\
 L_{K,p}(t)\le p(K-1)t}}
 \frac{Np+2pt+L_{K,p}(t)}{Np(1+t)}
 \right\}.}
\end{equation}
Only finitely many endpoints need checking: floor breakpoints and crossings
of the piecewise-linear constraints.
\end{theorem}
\begin{proof}
Write $q_{ij}=\Pp(Y_i=Y_j=1)$. The influence is $v_{ij}=(2q_{ij}-p)_+$,
including $v_{ii}=p$. The negative-correlation branch can contribute only
if $q_{ij}<(3p-1)/2$, impossible here.
Choose a maximal-gain root with row sum $p(1+t)$ and let its $K$ active
neighbors satisfy $q_{i0}>p/2$. If $t=0$ all rows equal $p$.
Otherwise let $Z$ count the ones in the active set.
Summing root joint probabilities gives conditional means $\mu_1,\mu_0$.
The convex envelope of $k(k-1)$ on integers is $\phi$, so
$\E[Z(Z-1)]\ge p\phi(\mu_1)+(1-p)\phi(\mu_0)$.
The ordered active-leaf influence sum is at least $L_{K,p}(t)$, and the full
sum is at least $Np+2pt+L_{K,p}(t)$.
Sum the $K$ inequalities saying each active leaf row is at most the root
row: its active-leaf influence sum is at most $p(K-1)t$.
Thus every law yields a feasible pair and the lower bound in \eqref{eq:rare}.

Conversely take a $\mathrm{Bernoulli}(p)$ root. Conditional on root $s$,
draw $Z$ on the two neighboring integers around $\mu_s$ with that mean,
then choose a uniform subset of $Z$ active leaves to label one.
Every active leaf has marginal $p$ and root influence $pt/K$.
Exchangeability makes its total ordered leaf influence exactly $L_{K,p}(t)$.
Take remaining labels independently $\mathrm{Bernoulli}(p)$, independent of
the block; their cross-influences vanish.
The feasibility constraint is exactly that the root row be maximal.
This attains every feasible displayed ratio. Between all specified
breakpoints the objective is linear fractional and is minimized at an
endpoint (or is constant).
\end{proof}

\begin{proposition}[Imbalance changes the asymptotic order]\label{prop:rareorder}
\pr\ For every fixed $0<p<1/2$, $\kappa_N(p)=\Theta(N^{-1})$.
Define $t_p=2p-1+\sqrt{2(1-p)(1-2p)}>0$. Then
\[
 \frac{3N-2}{N^2}\le\kappa_N(p)
 \le\frac{N+2(N-1)t_p}{N[1+(N-1)t_p]}.
\]
For $0<p\le1/3$,
$\lim_{N\to\infty}N\kappa_N(p)=2+t_p^{-1}$.
\end{proposition}
\begin{proof}
Homogeneous influences are symmetric, nonnegative, bounded by $p$.
If their maximal row is $D$, the total is at least $Np+2(D-p)$ and
$D\le Np$. Divide by $ND$.
For the upper bound use a root of probability $p$ and conditionally iid
leaves with probabilities
$a=p+\sqrt{(1-p)(1/2-p)}$ given root 1 and
$b=p(1-a)/(1-p)$ given root 0.
Leaves have marginal $p$; distinct leaf pairs have joint probability $p/2$
and zero influence. Root--leaf influence is $p(2a-1)=pt_p$.

For the limit, an optimizing pair in \eqref{eq:rare} has $t=\Omega(N)$,
$K\ge t=\Omega(N)$ and $L/p=O(N)$, by the upper bound and numerator $Np$.
The inequality $\phi(z)\ge z^2-z$ yields
\[
 L/p\ge[K^2 f_p(t/K)-K]_+,\qquad
 f_p(z)=\frac{z^2+2(1-2p)z-(1-2p)}{2(1-p)}.
\]
This increasing quadratic for $z\ge0$ has positive root $t_p$.
Thus $t/K\le t_p+o(1)$, so $t/N\le t_p+o(1)$.
Finally $N\kappa_N(p)\ge(N+2t)/(1+t)
=2+(N-2)/(1+t)$ gives the matching lower limit.
\end{proof}
\nc\ Exact rational endpoint enumeration for $N=136$, $p=3/34=12/136$
gives $\kappa_{136}(3/34)=2798/86615\simeq0.0323038735$.
An attaining law has $K=135$, $t=5004/91$, $L=0$,
$\mu_1=17289/182$, $\mu_0=21843/5642$, and leaf-pair probability $3/68=p/2$.
This fixes each marginal to $12/136$; it does not fix the realized count.
Under the uniform law on exactly 12 rare labels all query values coincide,
so the ratio is one. The finite LP can incorporate a fixed-count constraint.

\section{Different targets and noisy observations}
\begin{proposition}[Exact positivity condition]\label{prop:transfer}
\pr\ For arbitrary targets and binary observations, query $a$ has positive
Bayes Hamming value iff some positive-weight target satisfies
$|c_{ia}|>|m_i|$. If a set of targets of total weight at least $\omega$
has $|c_{ia}|-|m_i|\ge g>0$, then
\[
 \Delta(a)\ge\omega g/2,\qquad
 \Delta(a)/\Delta^\star\ge\omega g/(2R_0)\ge\omega g.
\]
\end{proposition}
\begin{proof}
Use the nonnegative summands in \eqref{eq:gain} and
$\Delta^\star\le R_0\le1/2$.
\end{proof}
This is a necessary-and-sufficient pointwise condition for positive value,
and a quantitative sufficient condition for a ratio. There is no unique
weakest named family of structural assumptions; any uniform ratio needs
quantitative transfer of this kind or an equivalent restriction.

\begin{example}[Disjoint targets]\label{ex:heldout}
\ce\ A fair target $T$ and a fair candidate $O_1$ can be independent, while
$O_2$ is a noisy copy of $T$ with $\E O_2\ne0$ and positive information value.
For example set $O_2=+1$ when $T=+1$ and, when $T=-1$, let $O_2=+1$ with
probability $1/2$. Then $\E O_2=1/2$, $\E TO_2=1/2$.
Margin uniquely selects $O_1$, with zero value, whereas $O_2$ has value $1/4$.
Hence target--pool separation admits no unconditional positive margin ratio.
\end{example}

\begin{theorem}[Common binary symmetric observation noise]\label{th:bsc}
\pr\ In the full-pool Hamming model observe $O_j=S_j\xi_j$, where the
one-query channel noise is independent of the world and
$\E\xi_j=\eta=1-2\epsilon\in[0,1]$, common to all candidates.
Latent-label margin still satisfies $\Delta(a)\ge\Delta^\star/(N-1)$.
For $\eta>0$ the constant is sharp. Observation-marginal margin has the same
ranking as latent margin.
\end{theorem}
\begin{proof}
Now $v_{ij}=(\eta|c_{ij}|-|m_i|)_+/2$ and
$v_{ii}=w_i=(\eta-|m_i|)_+/2=(u_i-\epsilon)_+$.
For every $i,j$, $0\le v_{ij}\le w_i$, and
\[
 w_i-v_{ij}=\tfrac12[\eta-\max\{|m_i|,\eta|c_{ij}|\}]_+.
\]
If $|m_a|\le|m_b|$, the right side at $(b,a)$ is no larger than at $(a,b)$.
Repeat the proof of Theorem \ref{th:margin} with $w_i$ in place of $u_i$.
The independent fair decoy and slightly biased duplicate cluster approach
the same sharp ratio. Finally $|\E O_j|=\eta|m_j|$.
\end{proof}
\begin{example}[Informative but heterogeneous noise]\label{ex:hetero}
\ce\ Take independent labels with probabilities $(0.10,0.09)$ and BSC
error rates $(0.20,0.01)$. Both latent and observation margin choose point 1.
Its Bayes value is zero, because its uncertainty $0.10$ is below its channel
error $0.20$. Query 2 has full-pool gain $(0.09-0.01)/2=0.04$.
Every channel is informative; that alone does not guarantee a positive ratio.
\end{example}
\begin{corollary}[Weighted overlapping targets]
\pr\ If the targets are the pool labels, possibly with known sign flips,
and all target weights lie in $[w_{\min},w_{\max}]$ with $w_{\min}>0$,
then noiseless or common-BSC margin has ratio at least
$w_{\min}/[w_{\max}(N-1)]$.
\end{corollary}
\begin{proof}
Sandwich each weighted sum of nonnegative influences between $w_{\min}$
and $w_{\max}$ times its unweighted sum.
\end{proof}

\section{Sequential budgets: proved bounds and unresolved sharp constant}
\begin{theorem}[Finite-horizon bounds]\label{th:budget}
\pr\ In the default noiseless model, sequential margin with exact updating
and $B$ distinct queries satisfies
\[
 G_{\rm margin}(B)\ge (B/N)R_0
 \ge (B/N)G_{\rm opt}(B),\quad 1\le B\le N,
\]
where the comparator is the optimal adaptive $B$-step policy.
For $B=1$ the sharper ratio is $1/(N-1)$; for $B=N$ it is one.
For $N\ge B+2$ the worst-case ratio is at most $B/(N-1)$.
For fixed $B$ and $N\to\infty$ these bounds establish the sharp leading
order $B/N$.
\end{theorem}
\begin{proof}
After $t$ queries, posterior Bayes risk is the sum of the $N-t$ remaining
uncertainties divided by $N$. Its largest uncertainty is at least
$NR_t/(N-t)$; observing that point reduces risk by at least its own
uncertainty divided by $N$. Thus
$\E[R_{t+1}\mid H_t]\le(1-1/(N-t))R_t$.
Iterating gives $\E R_B\le(N-B)R_0/N$.

For the upper construction take $B$ independent fair decoys and
$M=N-B$ copies of an independent $\mathrm{Bernoulli}(1/2+\eta)$ bit.
Margin spends all $B$ queries on decoys. For small $\eta$, an optimal policy
queries the cluster and $B-1$ decoys; further cluster queries are useless.
The ratio is
$B/[M(1-2\eta)+B-1]\to B/(N-1)$.
\end{proof}
\begin{example}[The tempting $B/(N-1)$ lower bound is false]\label{ex:budget}
\ce\ For $N=3$, $B=2$, let
\[
 P(010)=1/6,\quad P(100)=1/3-\varepsilon,\quad
 P(101)=1/3+\varepsilon,\quad P(110)=1/6,
 \qquad 0<\varepsilon<1/6.
\]
Uncertainties are $(1/6,1/3,1/3+\varepsilon)$, so margin first queries
coordinate 3. If it is zero, margin queries coordinate 2, leaving
unnormalized expected residual error $1/6$; if it is one the world is known.
Initial unnormalized risk is $5/6+\varepsilon$.
An optimal policy queries coordinate 2, then coordinate 3 if it is zero
and coordinate 1 otherwise, identifying the world exactly.
The ratio is $(4+6\varepsilon)/(5+6\varepsilon)\to4/5<1=B/(N-1)$.
\end{example}
\begin{theorem}[Sharp three-point, two-query constant]\label{th:n3b2}
\pr\ For $N=3$, $B=2$, the worst-case ratio of sequential margin to the
optimal adaptive policy is exactly $4/5$. This includes arbitrary random
tie-breaking within the margin maximizers.
\end{theorem}
\begin{proof}
We give a finite rational certificate proof; all certificate coefficients
are displayed below, so the proof does not rely on floating-point
optimization or a claimed exhaustive sample of priors.
Use unnormalized risks. By coordinate flips suppose all marginal
probabilities $p_i\le1/2$, and permute the first margin query to coordinate
3, so $p_3\ge p_1,p_2$. Write the eight world probabilities in lexicographic
order as $\pi=(\pi_{000},\ldots,\pi_{111})$.

For four masses in a $2$ by $2$ table, the smaller expected residual error
after revealing one of the two coordinates is the sum of the two smallest
masses. Indeed each row/column error selects two entries and is at least
that sum; at least one of the two partitions separates the two largest
entries, attaining it. Conditional on the first query, margin is optimal
on the remaining two coordinates by Theorem \ref{th:margin}.
Thus its residual risk $F_j$ after first querying $j$ is the sum of the two
smallest masses in each of the two $j$-slices.

For $j=1,2,3$, form a $36$ by $8$ matrix $L^{(j)}$ as follows.
List the indices of the four worlds with coordinate $j=0$ and those with
$j=1$, each in lexicographic order. List their six two-element subsets in
lexicographic order. For each pair of subsets, taking the zero-slice subset
as the outer loop, put a row with ones on their union and zeros elsewhere.
Then $F_j=\min_k L^{(j)}_k\pi$.
The initial risk is $R_0=\sum_i p_i$.
It suffices to prove $5F_3\le R_0+4F_2$; symmetry gives the comparison with
query 1 as well.

Let $z=(\pi,h)\ge0$ and $e=(1,1,1,1,1,1,1,1,0)$.
Define the $41$ by $9$ matrix $A$ by rows
\[
 A_k=(-L^{(3)}_k,1)\ (1\le k\le36),\quad
 A_{36+i}=(x_i,0)\ (1\le i\le3),\quad
 A_{40}=(x_1-x_3,0),\ A_{41}=(x_2-x_3,0),
\]
where $x_i$ is the length-eight vector of coordinate-$i$ labels in the
world order. Let $b_{37}=b_{38}=b_{39}=1/2$ and all other $b_i=0$.
At $h=F_3$ we have $Az\le b$ and $e^\top z=1$.
For each $k=1,\ldots,36$ set
$c_k=(4L^{(2)}_k+x_1+x_2+x_3,-5)$.
The table provides $\lambda_k$ and a nonnegative vector $v_k$ such that
\[
 c_k+A^\top v_k-\lambda_ke\ge0,\qquad
 d_k=\lambda_k-b^\top v_k\ge0.
\]
These are just nine rational inequalities and one rational equality per row,
directly checkable using the matrix definitions.
Therefore
$c_k^\top z\ge\lambda_ke^\top z-v_k^\top Az
\ge\lambda_k-v_k^\top b=d_k\ge0$.
Taking the minimum over $k$ proves $5F_3\le R_0+4F_2$.
Consequently $R_0-F_3\ge(4/5)(R_0-\min_jF_j)$.
The argument holds for every allowed first margin choice and its conditional
continuation, hence also their randomized mixtures.
Example \ref{ex:budget} with $\varepsilon\downarrow0$ proves sharpness.
\end{proof}

\small
\begin{longtable}{r r p{.65\textwidth} r}
\toprule $k$ & $\lambda_k$ & Nonzero entries of $v_k$ (index:value) & $d_k$\\
\midrule\endhead
1 & $1$ & $1:3,\ 30:1,\ 35:1,\ 40:1$ & $1$\\
2 & $2$ & $13:1,\ 15:1,\ 26:1,\ 27:1,\ 30:1,\ 37:1,\ 39:1$ & $1$\\
3 & $6/7$ & $3:12/7,\ 5:8/7,\ 6:2/7,\ 27:12/7,\ 36:1/7,\ 41:5/7$ & $6/7$\\
4 & $4/5$ & $13:7/5,\ 14:6/5,\ 17:3/5,\ 25:1/5,\ 35:8/5,\ 40:7/5$ & $4/5$\\
5 & $6/7$ & $6:9/14,\ 9:1/2,\ 15:13/14,\ 17:15/14,\ 35:13/7,\ 40:31/14,\ 41:1/2$ & $6/7$\\
6 & $2/3$ & $15:10/3,\ 29:1,\ 30:1/3,\ 35:1/3,\ 41:1$ & $2/3$\\
7 & $2$ & $5:1,\ 11:1,\ 22:2,\ 35:1,\ 37:1,\ 39:1$ & $1$\\
8 & $5/2$ & $1:3/4,\ 9:3/8,\ 10:3/8,\ 30:7/4,\ 32:3/8,\ 35:11/8,\ 39:3$ & $1$\\
9 & $3/2$ & $5:3/2,\ 9:1,\ 24:2,\ 35:1/2,\ 39:3/2$ & $3/4$\\
10 & $2$ & $13:1/3,\ 16:4/3,\ 17:1/3,\ 35:3,\ 38:1,\ 39:4/3$ & $5/6$\\
11 & $2/3$ & $5:1/3,\ 10:4/3,\ 11:5/3,\ 35:5/3,\ 40:1/3$ & $2/3$\\
12 & $2$ & $15:1/3,\ 18:5/3,\ 33:1/3,\ 35:8/3,\ 38:1,\ 39:5/3$ & $2/3$\\
13 & $2/3$ & $4:10/3,\ 28:1,\ 35:2/3,\ 40:1/3$ & $2/3$\\
14 & $2$ & $14:1/3,\ 16:4/3,\ 18:1/3,\ 30:3,\ 37:1,\ 39:4/3$ & $5/6$\\
15 & $2/3$ & $4:4/3,\ 6:5/3,\ 12:1/3,\ 30:5/3,\ 41:1/3$ & $2/3$\\
16 & $1/2$ & $16:2,\ 18:3/2,\ 28:1/2,\ 35:1,\ 40:1/2$ & $1/2$\\
17 & $0$ & $4:1,\ 10:1,\ 18:2,\ 35:1,\ 40:1$ & $0$\\
18 & $1/2$ & $18:7/2,\ 28:1,\ 35:1/2,\ 41:1/2$ & $1/2$\\
19 & $0$ & $19:1,\ 20:2,\ 35:2$ & $0$\\
20 & $0$ & $26:2,\ 27:1,\ 35:2$ & $0$\\
21 & $0$ & $21:3,\ 35:2$ & $0$\\
22 & $0$ & $19:1,\ 31:1,\ 35:3$ & $0$\\
23 & $0$ & $20:3/2,\ 35:7/2,\ 40:1,\ 41:1/2$ & $0$\\
24 & $0$ & $30:1,\ 32:1,\ 33:1,\ 35:2$ & $0$\\
25 & $0$ & $19:1/2,\ 26:3/2,\ 35:3,\ 40:5/2$ & $0$\\
26 & $0$ & $26:7/2,\ 35:3/2,\ 40:1/2$ & $0$\\
27 & $0$ & $24:1,\ 26:3,\ 35:1,\ 40:1,\ 41:1$ & $0$\\
28 & $0$ & $26:2,\ 35:3,\ 40:2,\ 41:1$ & $0$\\
29 & $0$ & $26:3/2,\ 35:7/2,\ 40:5/2,\ 41:1/2$ & $0$\\
30 & $0$ & $26:2,\ 30:3/2,\ 35:3/2,\ 40:1/2,\ 41:5/2$ & $0$\\
31 & $0$ & $20:1,\ 22:1,\ 24:1,\ 35:2$ & $0$\\
32 & $0$ & $26:1,\ 28:1,\ 30:2,\ 35:1$ & $0$\\
33 & $0$ & $21:1,\ 24:2,\ 35:2$ & $0$\\
34 & $0$ & $19:1,\ 34:1,\ 35:3$ & $0$\\
35 & $0$ & $22:2,\ 35:3,\ 41:1$ & $0$\\
36 & $0$ & $30:1,\ 32:1,\ 35:2,\ 36:1$ & $0$\\
\bottomrule
\end{longtable}
\normalsize
Every displayed certificate was verified in exact rational arithmetic.
The general finite-$N,B$ sharp ratio is still unresolved.
Random tie-breaking does not repair the unique-choice limiting construction.

\begin{proposition}[What randomization does establish]
\pr\ Uniform sampling without replacement achieves at least
$(B/N)R_0$ expected reduction. With a known joint law, an optimal adaptive
finite-horizon policy can be deterministic. In marginal-only one-step
selection randomization can strictly improve the worst-case guarantee,
as Theorem \ref{th:kappa} proves.
\end{proposition}
\begin{proof}
Each point is sampled with probability $B/N$ and its original error is then
zero; expected error at unqueried points cannot increase under Bayes
conditioning. For known-law optimality, backward induction maximizes a
linear function of the action probabilities at every history.
\end{proof}

\section{Misspecified selection and updating}
For acquisition-regret comparisons, $P$ is the reference posterior under the
correct generating law at the fixed data state. A truth-knowing oracle
conditions further on the realized world and is a different comparator.
The risk identity also allows a point-mass $P$, but then its Bayes-optimal
Hamming baseline is already zero.
\begin{theorem}[True value under a model-based update]\label{th:PQ}
\pr\ Let $P$ be the true joint law and $Q$ the learner's law.
For targets and observations as in Theorem \ref{th:binary}, write
$m_i^P,m_i^Q,c_{ij}^P,c_{ij}^Q$ for their moments.
The learner initially predicts $h_i=\sg(m_i^Q)$ and, after $O_j=s$,
predicts $g_{ij}(s)=\sg(m_i^Q+s c_{ij}^Q)$.
Define $\alpha_{ij}=(g_{ij}(+1)+g_{ij}(-1))/2$,
$\beta_{ij}=(g_{ij}(+1)-g_{ij}(-1))/2$.
Its true post-query risk and gain from its own pre-query risk are
\begin{align}
 F_{P,Q}(j)&=\tfrac12\sum_iw_i
 [1-\alpha_{ij}m_i^P-\beta_{ij}c_{ij}^P],\label{eq:PQrisk}\\
 U_{P,Q}(j)&=\tfrac12\sum_iw_i
 [(\alpha_{ij}-h_i)m_i^P+\beta_{ij}c_{ij}^P].\label{eq:PQgain}
\end{align}
These identities hold for any prescribed binary branch decisions $g$;
expressing $g$ through $Q$ moments assumes coherent conditioning under $Q$.
On a $Q$-null observation branch the displayed sign convention specifies
the learner's otherwise arbitrary action.
\end{theorem}
\begin{proof}
$g_{ij}(O_j)=\alpha_{ij}+\beta_{ij}O_j$ and
$\1_{\{g\ne T_i\}}=(1-T_ig)/2$. Take expectations under $P$ and subtract
the initial risk $\sum_iw_i(1-h_im_i^P)/2$.
\end{proof}
Away from ties, inactive pairs $|c_{ij}^Q|\le|m_i^Q|$ contribute zero gain.
For active pairs,
\[
 U_{P,Q}(j)=\tfrac12\sum_{i:\,|c_{ij}^Q|>|m_i^Q|}
 w_i[\sg(c_{ij}^Q)c_{ij}^P-\sg(m_i^Q)m_i^P].
\]
Thus the sum on the right being zero exactly characterizes zero true value;
it can also be negative. Positive $P$-optimal value is independently
characterized by $|c_{ij}^P|>|m_i^P|$ for some $i,j$ of positive weight.

\begin{theorem}[Pairwise discrepancy bounds]\label{th:regret}
\pr\ Define
\[
 E_j=\tfrac12\sum_iw_i(|m_i^P-m_i^Q|+|c_{ij}^P-c_{ij}^Q|),\qquad
 D_j=\sum_iw_i\max\{|m_i^P-m_i^Q|,|c_{ij}^P-c_{ij}^Q|\}.
\]
Let $j_Q$ maximize $\Delta_Q$ and $j_P$ maximize $\Delta_P$.
Then
\begin{align*}
 |U_{P,Q}(j)-\Delta_Q(j)|&\le E_j,&
 |\Delta_P(j)-\Delta_Q(j)|&\le E_j,\\
 \Delta_P(j_P)-U_{P,Q}(j_Q)&\le E_{j_P}+E_{j_Q},\\
 0\le F_{P,Q}(j_Q)-F_P(j_P)&\le(D_{j_Q}+D_{j_P})/2\le\max_jD_j.
\end{align*}
The terminal-risk comparison uses a common endpoint and is preferable to
comparing gains from two different baselines.
\end{theorem}
\begin{proof}
For the coherent sign rule, either the branch is constant and equals $h_i$
(zero gain), or $\alpha_{ij}=0$, $|\beta_{ij}|=1$, so
$|\alpha_{ij}-h_i|\le1$. Apply \eqref{eq:PQgain}; for $\Delta$ apply the
Lipschitz argument in Corollary \ref{cor:basic}. Maximization gives the gain
comparison. For terminal risks exactly one of $\alpha,\beta$ has magnitude
one, yielding $|F_{P,Q}(j)-F_Q(j)|\le D_j/2$.
Also $|F_P(j)-F_Q(j)|\le D_j/2$ because a maximum is 1-Lipschitz in the
sup norm. Insert $F_Q(j_Q)\le F_Q(j_P)$.
Nonnegativity follows from Bayes optimality under $P$.
\end{proof}
\begin{corollary}[Approximation and refit defect]
\pr\ If the estimated $Q$ scores have uniform error at most $\epsilon$,
the terminal bound acquires $2\epsilon$. If the actual branch predictor
$\widetilde g$ differs from the coherent predictor with weighted probability
\[
 \zeta_j=\sum_iw_i P_P(\widetilde g_{ij}(O_j)\ne g_{ij}(O_j)),
\]
add $\zeta_j$ for the selected candidate.
\end{corollary}
\begin{proof}
Score maximization loses at most $2\epsilon$.
Two predictions' zero-one losses differ by at most their disagreement
indicator. Sum and take expectations.
\end{proof}

\begin{example}[Selection failure with identical marginals]\label{ex:selection}
\ce\ Let all variables be fair. Under $P$, $O_1$ is independent of $T$ and
$O_2=T$; under $Q$, swap the two observations. The unique $Q$-optimal query
has zero true value, while the $P$-optimal query has value $1/2$.
All first moments agree exactly.
\end{example}
\begin{example}[Zero value even in the full pool]\label{ex:update}
\ce\ Take a fair root $R$ and two leaves. Under $P$ let the leaves be
$-RV_1,-RV_2$, with probabilities $1/2,1/4,1/4$ on
$(V_1,V_2)=(++),(+-),(-+)$.
Under $Q$ let them be $RV_1,RV_2$ with probabilities $3/5,1/5,1/5$.
All three marginals are fair. Under $P$ root--leaf moments are $-1/2$ and
the leaf-pair moment is zero. Under $Q$ they are $3/5,3/5,1/5$.
The unique $Q$-optimal query is the root (column sum $11/5$ versus $9/5$).
Its $Q$ update predicts all labels equal to the root. True posterior risk
is $(0+3/4+3/4)/3=1/2$, exactly the baseline: zero gain.
The $P$-optimal query is also the root, but its correct update gains $1/3$.
This isolates update failure from candidate misranking.
\end{example}
\nc\ Finite random-law checks for $N=2,\ldots,6$ (1,500 $P,Q$ pairs)
confirmed the direct-risk identities, terminal discrepancy bound, and
common-noise margin inequality. These checks supplement the proofs.
\nc\ A separate fixed-marginal check for $(.1,.1,.25)$ enumerated the six
vertices of its world-law polytope and 50,000 random convex mixtures.
The smallest tested ratio of the selector in Theorem \ref{th:twolevel}
was $20/23$, attained at a vertex. The proof, rather than this finite
sample, establishes the universal claim.

\section{Week 15 cross-audit and corrected propositions}
Inspected: \texttt{outputs/week15\_boundary\_acquisition/WEEK15\_REPORT.md},
\texttt{THEORY\_WEEK15.md}, \texttt{THEORY\_ERRATA.md}, the Week 14 theory,
and \texttt{src/week15\_bayes\_frontier.py}, \texttt{week15\_metric.py},
\texttt{week15\_metric\_confirm.py}. Empirical results were not rerun.
\small
\begin{longtable}{p{.07\textwidth}p{.13\textwidth}p{.70\textwidth}}
\toprule Item & Verdict & Minimal correction\\\midrule\endhead
B1 & Correct identity; gap &
XOR and orientation invariance are exact. Error-set perimeter is not a
displacement metric; delete its general identification with NSD/ASSD.\\
B2 & Gap &
Fixed-$r$, fixed-set iid-cloud consistency is correct. Fix the posterior
independently of the evaluation cloud or prove uniform convergence.
Require positive denominator. Fixed-$k$ kNN and shrinking-radius claims
need separate normalization/assumptions. Graph Dice is not NSD.\\
B3 & Wrong as exact; repairable &
An unspecified spread does not imply small overlap probability. State a
posterior law. For a uniform threshold posterior, normalized plug-in loss
is $1-r/(3s)$, not exactly one. Static risk does not prove the cause of
an acquisition run's collapse.\\
B4 & Correct theorem; empirical gap &
Fix state, posterior, metric, update, and coupling of potential labels.
Truth-information gap is $V_E-V_B$; $V_E-V_{\rm margin}$ also contains a
method gap. Separate-policy AULCs are not this decomposition.
The inspected implementation is Laplace/Monte Carlo NSD look-ahead.\\
M1 & Correct core; gaps &
For true weights and fixed radius/geometry, bounded-below measurable density
suffices. Smoothness at graph scale is unnecessary. Assume uniform density
consistency explicitly; the listed $k$ conditions do not establish it.
Fixed 95\% clipping changes the estimand. Radius and standardization also
vary between designs in the implementation.\\
E15-1 & Correct &
Exchangeable fixed counts, distinct queries, independent randomization,
both classes present; survival at zero is one.\\
E15-2 & Correct number; notation gap &
Use $(H_{108}^2+\sum_{i=1}^{108}i^{-2})/2=14.67254557\ldots$.
Standard $H_n^{(2)}$ alone usually means $\sum i^{-2}$.\\
E15-3 & Correct with convention &
Resolve ties and compute AUC on the same order. Half-credit AUC on unresolved
ties can violate the claimed relation to the query order.\\
E15-4 & Correct &
The convex dominance closure is contained in the implied set. State
consistency and strict-inequality conventions for the two feasibility LPs.\\
E15-5 & Correct &
The zero-budget survival term is one; the binomial sum starts at one.\\
\bottomrule
\end{longtable}
\normalsize

\begin{proposition}[Fixed-cloud continuum estimands]\label{prop:Ustat}
\pr\ Let $X_i$ be iid on bounded $\Omega$ with density $p\ge p_{\min}>0$,
and let $h$ be a bounded measurable symmetric pair kernel, fixed independently
of the cloud. Then, almost surely,
\[
 {1\over {n\choose2}}\sum_{i<j}\frac{h(X_i,X_j)}{p(X_i)p(X_j)}
 \longrightarrow\int_\Omega\int_\Omega h(x,y)\,dx\,dy.
\]
This applies to fixed-radius true, predicted, shared and mismatched cuts.
Ratios converge when their limiting denominators are positive.
A fixed posterior average independent of the cloud is allowed.
If $\|\widehat p-p\|_\infty\to0$ almost surely, plug-in weights have the
same limit.
\end{proposition}
\begin{proof}
The bounded weighted kernel obeys the U-statistic strong law; sampling
densities cancel. For posterior averaging use its integrated kernel,
conditioning first on fixed training data. Uniform density consistency
implies eventual $\widehat p\ge p_{\min}/2$ and uniformly converging inverse
products. Use continuous mapping for the ratio.
\end{proof}
The unweighted uniform-cloud limit has factor $|\Omega|^{-2}$.
Density continuity is unnecessary for true weights.
If clipping thresholds converge to a finite $c$, the limit is instead
$\int\!\int h(x,y)\min\{1,c\,p(x)p(y)\}\,dx\,dy$ by dominated convergence.
Fixed quantile clipping therefore generally retains density dependence.

\begin{proposition}[Exact flat-boundary repair]\label{prop:flat}
\pr\ On $(-L,L)$ let $A_\theta=\{x>\theta\}$, with all thresholds farther
than $r$ from the endpoints. For ordered pairs with $|x-y|<r$,
\[
 J_r(A_\theta)=r^2,\qquad
 S_r(A_\theta,A_0)=(r-|\theta|)_+^2.
\]
If $\Theta\sim\mathrm{Uniform}[-s,s]$, $s\ge r$, $L>s+r$, then
\[
 \E S_r=r^3/(3s),\quad
 \E J_r(A_\Theta\triangle A_0)=2r^2(1-r/(3s)).
\]
A constant predictor has mismatch $r^2$, strictly smaller.
Normalized ratio-of-expectations losses are $1-r/(3s)$ and one, respectively.
\end{proposition}
\begin{proof}
Crossing pairs form two triangles of total area $r^2$; shared crossings
form triangles with side $(r-|\theta|)_+$. Integrate the square against
the uniform density and use mismatch = true + predicted $-2$ shared.
\end{proof}
\ce\ For $\theta=r/2$, graph Dice is $1/4$ while NSD with tolerance $r$
is one. These functionals differ.
\pr\ If expected true cut mass is positive, normalized static loss is
at most one for every predictor because shared mass is nonnegative.
A constant predictor equals one, so it is never strictly preferred over
a predictor with positive shared mass. This does not prove NSD improvement
by the resulting acquisition.

\begin{proposition}[Oracle decomposition]\label{prop:oracle}
\pr\ Fix data, a finite common action set, a joint law of truth $f$ and all
potential labels $(Y_c)$, and an integrable gain $G(c,Y_c,f)$ for a fixed
update procedure. Then
\[
 V_R=\E\max_cG(c,Y_c,f)\ \ge\
 V_E=\E_f\max_c\E[G(c,Y_c,f)\mid f]\ \ge\
 V_B=\max_c\E G(c,Y_c,f).
\]
$V_B$ dominates any rule independent of unknown truth and potential labels
conditional on the data.
\end{proposition}
\begin{proof}
Apply convexity of the maximum to conditional expectation twice.
A non-anticipating randomized rule averages fixed action values.
\end{proof}
Label marginals alone do not define $V_R$; a coupling is required.
The chain holds for any such coupling. Values need not be nonnegative for
a fixed non-Bayes prediction/update procedure evaluated by NSD.

\nc\ The inspected frontier code uses \texttt{S=24}, Laplace Gaussian
posterior samples, and fresh fantasy-label Laplace fits. A well-specified
prior does not make these exact posterior updates.
The reported $0.0155/0.0645\simeq24\%$ compares absolute gains.
Relative to reported margin gain $-0.007$, captured headroom is
$(0.0155+0.007)/(0.0645+0.007)\simeq31.5\%$.
Both descriptions must name their denominator.

\begin{proposition}[Precise errata statements]\label{prop:errata}
\pr\ For $n$ iid points with independent continuous coordinates in three
dimensions, expected minimal Pareto-front size is
$\sum_{i=1}^n H_i/i=(H_n^2+\sum_{i=1}^ni^{-2})/2$.
In a strict score order with $r$ rare and $m$ common labels, let $a_i$
count common labels preceding rare label $i$, in rare-label order.
First rare position is $a_1+1$ and
$m(1-\mathrm{AUC})=r^{-1}\sum_i a_i\ge a_1$.
Equality holds iff rare labels are consecutive.
\end{proposition}
\begin{proof}
Sort in the first coordinate. Point $i$ is minimal iff its two other
coordinates are undominated among the first $i$ points. Their expected
minimal count is $H_i$, by record probabilities, so exchangeability gives
probability $H_i/i$. Sum and rearrange the harmonic double sum.
For AUC each rare--common inversion contributes one pair.
The nondecreasing $a_i$ have average equal to their minimum iff all equal.
\end{proof}
\ce\ Positive point $(1,0)$ and negative point $(0,1)$ under positive
linear weights imply $w_1>w_2>0$, $t\in[w_2,w_1)$.
Point $(2,-1/2)$ must be positive but is outside the upward closure of
the positive convex hull, proving E15-4's containment can be strict.
\pr\ Exact implication is two candidate-label feasibility checks.
For finitely many constraints, strictly positive weights and positive-class
margins can be jointly scaled to at least one; negative scores remain
at most zero. These are LP constraints. Exclude inconsistent training data.

\section{Round-1 results retained: prediction, discovery, and stopping}
\begin{example}[Better marginals, worse acquisition]\label{ex:prediction}
\ce\ Use a fair decoy and $M=N-1$ copies of
$\mathrm{Bernoulli}(1/2+\eta)$, $N\ge3$.
Model A reports true probabilities. Model B reports $1/2+2\eta$ on the
decoy, $1/2-\eta/2$ on one cluster point and $1/2+3\eta/2$ on the others,
for small $\eta>0$.
A is strictly better in expected Brier/log loss at every point and has
classification error lower by $2\eta/N$.
B's average Brier excess is $\eta^2(M+24)/(4(M+1))$ and maximum probability
error $2\eta$. Yet A's margin queries the decoy and B's queries the cluster.
Under a common ideal update their acquisition-value ratio tends to $1/M$.
\end{example}
\begin{proof}
Proper scores minimize uniquely at the true probabilities. B changes the
majority decision on one cluster point only. Its squared errors sum to
$4\eta^2+(3\eta/2)^2+(M-1)(\eta/2)^2$.
Its smallest absolute margin is the misclassified cluster point.
\end{proof}

\begin{theorem}[Finite-population discovery]\label{th:discovery}
\pr\ In a uniform random order of $n_0,n_1\ge1$ labels, $N=n_0+n_1$,
let $T$ be first both-class discovery. Then
\[
 \Pp(T>0)=1,\quad
 \Pp(T>k)=\frac{\binom{n_0}k+\binom{n_1}k}{\binom Nk}\ (k\ge1),\quad
 \E T=\frac{N+1}{n_0+1}+\frac{N+1}{n_1+1}-1.
\]
For first discovery $D$ of a designated class with $r$ members,
\[
 \Pp(D>k)=\frac{\binom{N-r}k}{\binom Nk}
 =\frac{\binom{N-k}r}{\binom Nr}
 \le\min\{(1-r/N)^k,(1-k/N)^r\},\quad
 \E D=\frac{N+1}{r+1}.
\]
The geometric mean $N/r$ exceeds this by $(N-r)/(r(r+1))$.
\end{theorem}
\begin{proof}
Survival means a monochromatic sample; its two events are disjoint for
$k\ge1$. If $D_0,D_1$ are first class positions, $\min(D_0,D_1)=1$,
hence $T=D_0+D_1-1$. The minimum of $r$ uniform distinct positions has
expectation $(N+1)/(r+1)$, by exchangeability of the $r+1$ gaps.
The two survival product forms give the bounds.
\end{proof}
\begin{corollary}[Exchangeable minimax and limits]
\pr\ Under the uniform fixed-count assignment prior, all adaptive distinct
query rules with truth-independent randomization have this discovery law.
Uniform query ordering is minimax over assignments for expected discovery
time. For fixed $r$, $T/N\Rightarrow\mathrm{Beta}(1,r)$.
If $r\to\infty$, $r/N\to0$, then $(r/N)T\Rightarrow\mathrm{Exp}(1)$.
\end{corollary}
\begin{proof}
Unqueried labels remain exchangeable after every history; the transition
probabilities depend only on remaining counts. A policy's worst assignment
is at least its prior average, and uniform ordering equalizes assignments.
The rare-class survival products tend to $(1-x)^r$ and $e^{-x}$ at the
respective scalings. Except with probability $r/N$, the first label is
common and $T$ equals rare-class discovery time.
\end{proof}
\nc\ At $N=136,r=12$: $\E T=10.6344615\ldots$,
$P(D>9)=0.42394419\ldots$, $P(D>16)=0.2078709166\ldots$;
rare-class 95\% and 99\% discovery quantiles are 29 and 42.
These are design calculations, not a frozen thesis endpoint estimate.

\begin{proposition}[Exact enrichment condition]\label{prop:tail}
\pr\ A subset of size $M$ containing $s\ge1$ designated-class members has
first-discovery time stochastically no larger than a full pool of size $N$
containing $r$, iff $s/M\ge r/N$.
Both-class startup additionally needs an opposite-class anchor or a separate
argument.
\end{proposition}
\begin{proof}
Necessity is the first-draw probability. For sufficiency compare hazards
$s/(M-j)$ and $r/(N-j)$ while survival is possible:
$sN-rM+j(r-s)\ge0$, because the subset has $s\le r$.
Multiply the complementary hazards; subset exhaustion sets survival to zero.
\end{proof}
\begin{example}[Weak global prediction can aid startup]\label{ex:startup}
\ce\ Let $N=r^2$, $r\ge4$. A three-point high-score subset contains one
rare and two common labels; outside it are $r-1$ rare labels.
Both conditional rare probabilities are below $1/2$, so the Bayes classifier
remains constant. Rare-oriented AUC is
$1/2+\frac12(1/r-2/(N-r))\to1/2$.
Calibrated Brier improvement is $(N-3r)^2/[3N^2(N-3)]\to0$ and log-loss
gain is at most $h_2(3/N)\to0$.
Both-class discovery within the subset has mean $7/3$; full-pool mean is
asymptotic to $r$.
\end{example}
\begin{proof}
Count binary-score pairs for AUC; compute the variance of the conditional
rare-class probability for Brier improvement. Mutual information is bounded
by the subset indicator entropy. Apply the discovery formula to each pool.
\end{proof}
\begin{example}[Same startup time, different information]\label{ex:history}
\ce\ For fair $\Theta$, let labels be $(0,1,\Theta,1-\Theta)$.
Querying the first two or last two gives $T=2$ in both cases, but posterior
full-pool risk is respectively $1/4$ or zero.
\end{example}
\begin{theorem}[Stopping-time information]\label{th:stop}
\pr\ Let $\Theta$ be a finite identifiable family of deterministic binary
worlds, each containing both classes. Let $U$ be an independent policy
seed, $H_T$ the history at first both-class discovery, and $Q\ge T$ a
zero-error identification stopping time. In bits,
\[
 I(\Theta;H_T\mid U)=H(Y_{\rm first},T\mid U)\le1+H(T),\qquad
 \E Q\ge\E T+H(\Theta)-1-H(T).
\]
If $T\le b$, the information is at most $\log_2[2(b-1)]$;
if $T$ is deterministic it is at most one bit.
\end{theorem}
\begin{proof}
Given $U$, all labels before $T$ equal the first label and the last label
is opposite. Thus $(Y_{\rm first},T)$ reconstructs all queries and labels;
conversely the history gives this pair. History is deterministic given
$(\Theta,U)$, proving the information identity.
Conditional on $(H_T,U)$ the remaining zero-error decision tree is a binary
prefix code, with mean length at least $H(\Theta\mid H_T,U)$.
Average and use independence of $U$. There are at most $2(b-1)$ pairs.
\end{proof}

\section{Round-1 results retained: boundary metrics and structural priors}
\begin{example}[A nearest-opposite-label band can reverse boundary ranking]
\label{ex:q20}
\ce\ Let $m\ge4$, $0<\epsilon<1/(5m)$, and use the $5m$ test points
$-2,-\epsilon,\epsilon,2\epsilon,\ldots,(5m-2)\epsilon$ with truth
$y(x)=\1_{\{x>0\}}$. Let $q20$ be the $m$ points of smallest distance to
an opposite-labelled test point. Then
$q20=\{-\epsilon,\epsilon,\ldots,(m-1)\epsilon\}$, with a strict cutoff.
Threshold A at $-1$ has band accuracy $(m-1)/m$ and boundary error one.
Threshold B at $(m-3/2)\epsilon$ has band accuracy $2/m$ and boundary
error $(m-3/2)\epsilon$. Both have nonempty boundaries; the band has both
classes. Taking $\epsilon=1/(100m^2)$ makes the geometric reversal arbitrarily
large. Embed the construction in four dimensions by constant other
coordinates; a common affine scaling of the varying coordinate preserves it.
\end{example}
\begin{proof}
Nearest-opposite distances are $2+\epsilon$ at $-2$, $2\epsilon$ at
$-\epsilon$, and $(j+1)\epsilon$ at $j\epsilon$.
They give the stated strict band cutoff. A misclassifies only its negative
band point. B correctly classifies that point and the largest positive band
point. Distances of the thresholds from zero give the geometric errors.
Training anchors at $\pm R$, $R>4$, can be included without changing this
test-band construction.
\end{proof}

\begin{proposition}[Graph cuts and localization]\label{prop:cut}
\pr\ On a graph $G=(V,E)$ let $C_y=\{\{i,j\}\in E:y_i\ne y_j\}$ and
$A=\{i:y_i\ne\widehat y_i\}$. Then
$C_y\triangle C_{\widehat y}=\delta_G(A)$.
Zero cut mismatch means the two labelings differ by an independent global
flip on each connected component; one known label per component removes
this ambiguity. Also
$|C_y\triangle C_{\widehat y}|\le d_{\max}|A|$.
Cut disagreement count alone cannot measure displacement.
\end{proposition}
\begin{proof}
For each edge,
$(y_i\oplus y_j)\oplus(\widehat y_i\oplus\widehat y_j)
 =(y_i\oplus\widehat y_i)\oplus(y_j\oplus\widehat y_j)$.
An error set with no crossing edge is a union of connected components.
Counting incident edges gives the degree bound. On a path, any two different
single-cut labelings disagree on exactly two cut edges, irrespective of
their separation.
\end{proof}

\begin{proposition}[A finite-pool boundary-location metric]\label{prop:metric}
\pr\ For geometric edges define
\[
 \rho(\{x_1,x_2\},\{y_1,y_2\})
 =\min_{\pi\in\mathfrak S_2}\max_{i=1,2}\|x_i-y_{\pi(i)}\|.
\]
Hausdorff distance induced by $\rho$ is a metric on nonempty finite edge
sets. Add the empty set by assigning its distance to every nonempty set a
constant $D_0\ge\operatorname{diam}(\Omega)>0$.
Applied to cut sets it is a metric on labelings modulo component flips;
anchors give identity of labelings. On an equally spaced path, two single
cuts $i,j$ have distance $|i-j|h$.
\end{proposition}
\begin{proof}
The optimal endpoint matchings compose, proving the triangle inequality
for $\rho$; the other metric properties follow from endpoint matching.
Hausdorff distance inherits these properties on finite sets.
All nonempty distances are at most $\operatorname{diam}(\Omega)$, so adding
the empty set preserves the triangle inequality. Use
Proposition \ref{prop:cut} for the quotient. For path edges, matching left
to left and right to right minimizes the maximal displacement.
\end{proof}

\begin{theorem}[Geometric approximation by graph cuts]\label{th:geometry}
\pr\ Let $\Omega$ be bounded and convex and $X$ an $\epsilon$-net.
Connect all pairs at distance at most $6\epsilon$.
Let $A,B$ have compact smooth boundaries with two-sided tubular neighborhoods
of radius strictly greater than $2\epsilon$, contained in $\Omega$.
Label $X$ by membership in each set. For their cut sets,
\[
 \left|H_\rho(C_A,C_B)-H(\partial A,\partial B)\right|\le12\epsilon.
\]
\end{theorem}
\begin{proof}
At each boundary point take normal offsets $\pm2\epsilon$ and net points
within $\epsilon$ of them. The tubular condition puts these net points on
opposite sides; their distance is at most $6\epsilon$ and their edge
midpoint is within $3\epsilon$ of the boundary point.
Conversely every crossing edge meets the boundary along its segment
(convexity keeps the segment in $\Omega$), so its midpoint is within
$3\epsilon$ of a boundary point.
Each midpoint set is therefore within Hausdorff distance $3\epsilon$ of
its boundary. Their mutual Hausdorff distance differs from the true one
by at most $6\epsilon$.
For two edges of length at most $6\epsilon$, midpoint distance is at most
$\rho$, which is at most midpoint distance plus $6\epsilon$.
The same comparison holds for Hausdorff distances, giving $12\epsilon$.
\end{proof}

\begin{proposition}[Moment hierarchy]\label{prop:hierarchy}
\pr\ Current Bayes vertex Hamming risk uses first moments; one-query vertex
Hamming value uses second moments. Current cut-Hamming risk for the
edgewise Bayes cut predictor uses second label moments. One-query
cut-Hamming value can require third moments.
\end{proposition}
\begin{proof}
The first two claims follow from Theorem \ref{th:binary}.
The spin encoding a cut on $\{i,j\}$ is $-S_iS_j$, so its mean is a second
moment and its cross-moment with a queried $S_k$ is a third moment.
For strict necessity compare three independent fair bits with the uniform
law on $000,011,101,110$. All first and second moments agree.
On the triangle the current edgewise Bayes cut risk is $3/2$ for both.
After one bit, independent bits give zero reduction, whereas even parity
determines the opposite edge and gives reduction $1/2$.
The edgewise predictor need not itself be realizable as a vertex cut;
this distinction is part of the stated loss/action model.
\end{proof}

\begin{proposition}[A physics prior can help one target distribution and hurt another]
\label{prop:physics}
\pr\ At two independent design sites minimize
$\sum_i n_i(g_i-f_i)^2+\lambda\sum_i(g_i-\mu_i)^2$.
Then $g_i-f_i=\lambda(\mu_i-f_i)/(n_i+\lambda)$.
Compared with the zero prior, weighted squared-risk change is
\[
 \sum_i q_i\left(\frac{\lambda}{n_i+\lambda}\right)^2
 (\mu_i^2-2\mu_if_i).
\]
For $f=(10,1)$, $\mu=(10,3)$, $n=(9,1)$, $\lambda=1$,
target weights $(.9,.1)$ give risks $.4,.1,.925$ for the physics prior alone,
physics plus correction, and zero-prior correction.
Weights $(.1,.9)$ give $3.6,.9,.325$ respectively.
\end{proposition}
\begin{proof}
Differentiate the strictly convex objective coordinatewise, substitute the
solution, and expand the squared-error difference.
The stated risks follow by direct substitution.
\end{proof}
\begin{proposition}[Smoothness alone does not exclude finite sign worlds]
\label{prop:smooth}
\pr\ On distinct finite locations with minimum spacing $d_{\min}$, every
binary sign assignment is realized by an $L$-Lipschitz function, for every
$L>0$. It is also realized by a function of arbitrarily small positive
norm bound in an RKHS whose pool Gram matrix is strictly positive definite.
\end{proposition}
\begin{proof}
Assign values $\alpha s_i$ with $0<2\alpha\le Ld_{\min}$ and extend by
the McShane formula $\min_i(\alpha s_i+L\|x-x_i\|)$.
For an RKHS with Gram matrix $K$, the interpolant with coefficient vector
$K^{-1}\alpha s$ has squared norm $\alpha^2s^\top K^{-1}s$, which can be
made arbitrarily small.
\end{proof}
Known amplitude exceeding $L$ times a coverage radius can certify a sign;
binary signs alone do not supply that amplitude. These are elementary
extension/interpolation facts, not novelty claims.

\section{Concrete predictions for a Laplace GP classifier}
Let $Q_f=\mathcal N(\mu,\Sigma)$ be the current Laplace latent approximation.
Two target definitions must be separated. For \emph{persistent observed labels},
draw conditionally independent $Y_i$ once given $f$, with
$P(Y_i=1\mid f)=\sigma(f_i)$, and reveal that same $Y_j$ on query.
Then for distinct indices,
\[
 m_i=\E_{Q_f}\tanh(f_i/2),\qquad
 c_{ij}=\E_{Q_f}\tanh(f_i/2)\tanh(f_j/2),\qquad c_{jj}=1.
\]
For \emph{latent-sign targets} $T_i=\operatorname{sign}(f_i)$ and a logistic
observation $O_j$, instead use
\[
 m_i=\E_{Q_f}\operatorname{sign}(f_i),\qquad
 c_{ij}=\E_{Q_f}\operatorname{sign}(f_i)\tanh(f_j/2).
\]
In particular $c_{jj}=\E|\tanh(f_j/2)|<1$ for finite nonsingular latents.
These moments are computable by one- or two-dimensional Gaussian quadrature.
Coherent fantasy conditioning reweights $Q_f$ by the observation likelihood;
a fresh Laplace fit generally gives a different law.

Each row below gives a falsifiable calculation, not an assertion that a GP
can realize every extremal binary law or that Hamming rankings predict NSD.
\small
\begin{longtable}{p{.24\textwidth}p{.68\textwidth}}
\toprule Result & Computable prediction / next check\\\midrule\endhead
Binary identity, Theorem \ref{th:binary}&
For every candidate, quadrature moments and direct coherent two-branch
posterior Hamming risks must agree. Test persistent-label and latent-sign
targets separately; substituting latent covariance for $c_{ij}$ should not
be expected to agree.\\
Sharp margin and balanced $\kappa_N$&
For persistent-label targets and coherent updating, verify
$\Delta_{\rm margin}/\Delta^\star\ge1/(N-1)$.
At zero latent means, all marginals are fair; verify uniform ratio
$\ge\kappa_N$. Failure beyond numerical tolerance diagnoses inconsistent
moments, targets, or updates. Attainment by a particular GP kernel is not
predicted.\\
General/two-level marginals, Theorem \ref{th:LP}&
On a small pool, compute moments and the finite LP using only the marginal
vector. Its worst-law guarantee must hold for the GP joint law.
At two points with unequal uncertainties, margin must be Hamming-optimal
under coherent revelation.\\
Exact two-level family, Theorem \ref{th:twolevel}&
Tune a three-point Gaussian latent law's marginal means to persistent-label
probabilities $(.1,.1,.25)$. Query with probabilities
$(5/46,5/46,18/23)$; the expected coherent Hamming gain divided by the best
candidate gain must be at least $20/23$. A GP-restricted optimum may be
better than this universal-law guarantee.\\
Rare marginals, Theorem \ref{th:rare}&
Use a constant-mean equal-variance latent prior tuned so every label marginal
equals $3/34$. For $N=136$, its uniform-query ratio cannot fall below
$2798/86615$. Compare with the constant-count posterior separately.
The value is a universal-law lower bound, not a claim that this GP attains it.\\
Positive transfer, Proposition \ref{prop:transfer}&
For latent-sign targets and logistic observations, mark candidates with
no pair satisfying $|c_{ij}|>|m_i|$. Their exact one-step Bayes Hamming
gain must be zero despite positive mutual information or observation entropy.\\
Common/heterogeneous noise, Theorem \ref{th:bsc}&
Add a common independent BSC channel to persistent labels and verify the
margin bound using $\eta c_{ij}$. With candidate-dependent noise, search
for margin candidates with no decision-crossing pair; positive information
alone will not prevent zero gain.\\
Budgets, Theorem \ref{th:budget}&
On a small finite label law induced by the GP, condition exactly and compare
margin with dynamic programming for $B=2,3$. Verify the $B/N$ bound.
For $N=3,B=2$, verify the stronger sharp $4/5$ bound from
Theorem \ref{th:n3b2}. Refitting a Laplace model is a separate experiment
and need not obey these bounds.\\
Misspecification, Theorems \ref{th:PQ}--\ref{th:regret}&
In a synthetic GP world with a controlled reference posterior $P$, compare
direct branch risks with the $P,Q$ moment formula; terminal regret must
be at most $\max_jD_j$. Identical marginal calibration does not remove the
$c^P-c^Q$ term. A single realized ground-truth function is not itself this
posterior $P$.\\
Update defect&
At fixed states, compare coherent likelihood reweighting of the current
Laplace Gaussian with fresh Laplace fantasy fits. Record weighted
decision disagreement $\zeta_j$ and verify its risk bound.
Repeat numerical integration/Monte Carlo precision checks before ranking
the 24-sample frontier estimates.\\
B1--B3, graph versus displacement&
For GP posterior draws near a flat interface, vary posterior location spread
and graph radius separately. The uniform-threshold control must recover
$1-r/(3s)$ normalized loss. Plot graph Dice and NSD on the same displaced
thresholds; the analytic discrepancy at displacement $r/2$ is $1/4$ versus one.\\
B4, oracle decomposition&
At the same frozen states, use a common coupled fantasy-label array and
the same refit/metric rule for all oracle levels. Population estimates
should satisfy the Jensen chain as sampling error shrinks; separate
policy AULCs need not form an exact decomposition.\\
M1, Proposition \ref{prop:Ustat}&
Freeze scaling and radius; with known densities and no clipping, replicate
cloud sampling and check convergence to one common pair-integral Dice.
Then change clipping, density estimation, and radius one at a time.
Fixed clipping is predicted to retain density dependence.\\
Errata / discovery&
Under exchangeable fixed-count posterior worlds, feature-adaptive queries
cannot improve discovery over random. Verify $t=0$ separately and reproduce
the $14.6725$ front calculation as an arithmetic control, not as evidence
for a GP acquisition improvement.\\
\bottomrule
\end{longtable}
\normalsize

\section{What remains unresolved}
\cj\ The exact general finite-$N,B$ sequential-margin competitive ratio
remains open here. The $N=3,B=2$ case is now proved to be $4/5$
by the rational certificates in Theorem \ref{th:n3b2}.
The general-marginal LP is exact but may be computationally prohibitive;
the compact rare-marginal formula currently covers homogeneous
$p\le1/3$ (and $p\ge2/3$ by flipping labels).
There is no claim of a sharp GP-restricted minimax constant, a universal NSD
margin ratio, or a publication-priority theorem.
The user's independent round-1 checks and the scratch numerical checks
support specified formulas only; they do not resolve these open questions.
The requested Astra-ultra subagents encountered their usage limit; the
extensions assembled here did not receive a completed independent second
proof audit. Source environments and references passed static checks.
The native LaTeX compiler returned a platform-directory error
(\texttt{Unable to find standard directories for platform}); PDF compilation
and visual layout therefore remain unverified.

\begin{thebibliography}{99}
\bibitem{ml18icml} S. Mussmann and P. Liang (2018).
On the Relationship between Data Efficiency and Error for Uncertainty Sampling.
\url{https://proceedings.mlr.press/v80/mussmann18a/mussmann18a.pdf}
\bibitem{ml18nips} S. Mussmann and P. Liang (2018).
Uncertainty Sampling is Preconditioned Stochastic Gradient Descent on Zero-One Loss.
\url{https://arxiv.org/abs/1812.01815}
\bibitem{eer} S. Mussmann et al. (2022).
Active Learning with Expected Error Reduction.
\url{https://arxiv.org/abs/2211.09283}
\bibitem{liu} S. Liu and X. Li (2023), original version.
Understanding Uncertainty Sampling (later titled Understanding Uncertainty
Sampling via Equivalent Loss).
\url{https://arxiv.org/pdf/2307.02719v1}
\bibitem{epig} F. Bickford Smith et al. (2023).
Prediction-Oriented Bayesian Active Learning.
\url{https://proceedings.mlr.press/v206/bickfordsmith23a/bickfordsmith23a.pdf}
\bibitem{tal} J. H\"ubotter et al. (2024). Transductive Active Learning:
Theory and Applications. Conference version.
\url{https://openreview.net/pdf/a4d2fe62468dee7d8a39e169ecd8c2cfd078a896.pdf}
\bibitem{wang} C. Wang, S. Sun and R. Grosse (2021).
Beyond Marginal Uncertainty: How Accurately can Bayesian Regression Models
Estimate Posterior Predictive Correlations?
\url{https://proceedings.mlr.press/v130/wang21g/wang21g.pdf}
\bibitem{pass} B. Pass and S. Spektor (2017 preprint).
On Khintchine type inequalities for pairwise independent Rademacher random variables.
\url{https://arxiv.org/abs/1708.08775}
\bibitem{gap} L. Galli, K. Kaparis and A. N. Letchford (2012).
Complexity results for the gap inequalities for the max-cut problem.
Includes discussion of earlier Laurent--Poljak gap inequalities.
\url{https://www.lancaster.ac.uk/staff/letchfoa/articles/2012-gap-ineqs-complexity.pdf}
\bibitem{update} F. Bickford Smith, A. Foster and T. Rainforth (2024).
Making Better Use of Unlabelled Data in Bayesian Active Learning.
\url{https://proceedings.mlr.press/v238/bickford-smith24a/bickford-smith24a.pdf}
\bibitem{batch} K. Hu and S. Mussmann (2026, arXiv v3).
Batch Bayesian Active Learning with Partial Batch Label Sampling.
\url{https://arxiv.org/pdf/2510.09877}
\end{thebibliography}
\end{document}

